\documentclass[10pt]{article}
\usepackage{tmlr}
\usepackage{amsmath,amssymb,amsthm,mathtools}
\usepackage{microtype,booktabs,array}
\usepackage{hyperref}
\usepackage{url}
\hypersetup{pdftitle={Sharp Class-Count Bounds for Calibration Reuse in Temperature-Scaled Conformal Prediction},pdfauthor={},pdfsubject={Statistical learning theory and conformal prediction}}
\providecommand{\doi}[1]{\href{https://doi.org/#1}{doi: #1}}
\newtheorem{theorem}{Theorem}[section]
\newtheorem{lemma}[theorem]{Lemma}
\newtheorem{proposition}[theorem]{Proposition}
\newtheorem{corollary}[theorem]{Corollary}
\theoremstyle{definition}\newtheorem{definition}[theorem]{Definition}
\theoremstyle{remark}\newtheorem{remark}[theorem]{Remark}
\newcommand{\R}{\mathbb R}
\newcommand{\E}{\mathbb E}
\newcommand{\Prob}{\mathbb P}
\newcommand{\1}{\mathbf 1}
\newcommand{\HH}{\mathcal H}
\newcommand{\VC}{\operatorname{VC}}
\newcommand{\Cov}{\operatorname{Cov}}
\newcommand{\Beta}{\operatorname{Beta}}
\newcommand{\risk}{\mathfrak R}
\newcommand{\pos}[1]{\left[#1\right]_+}
\DeclareMathOperator*{\argmin}{arg\,min}
\title{Sharp Class-Count Bounds for Calibration Reuse\\in Temperature-Scaled Conformal Prediction}
\author{Anonymous authors}
\date{}
\begin{document}
\maketitle
\begin{abstract}
How much coverage can be lost when a single temperature is chosen using the same observations that calibrate a conformal prediction set? For nonrandomized least-ambiguous classification (LAC) with a fixed $K$-class softmax model, we prove that the temperature-and-threshold class has growth function at most $m+1+(K-2)m(m-1)$ and VC dimension $(1+o(1))\log_2 K$. The key lower construction is a local realization theorem: one classifier, fixed before calibration, can expose $K-1$ independent score coordinates at prescribed temperatures. This gives an exact Beta-order-statistic law for an adversarial selector's coverage. Together with a continuum upper bound, it yields a uniform finite-sample worst-selector marginal deficit of order $\min\{1,\sqrt{\log K/n}\}$ at interior finite-rank coverage levels, including the saturation regime. The lower bound persists under bounded logits and any nonempty open positive-temperature interval, although no uniform numerical margin is guaranteed. A three-class example illustrates the construction and retains a valid likelihood-tuning control. The results characterize unmodified same-data LAC; they neither establish failure of ordinary likelihood tuning nor exclude validity-preserving reuse methods.
\end{abstract}

\section{Introduction}
Temperature scaling changes a classifier's softmax probabilities through one global parameter \citep{guo2017}. In conformal classification, these probabilities also determine a score whose empirical quantile defines a prediction set. A fixed score admits the usual rank-based coverage guarantee. Choosing its temperature on the calibration observations makes the score data-dependent, so that argument no longer applies without further structure.

The general possibility of selection bias is well understood. The unresolved structural question addressed here is more specific: \emph{what coverage cost can the actual one-temperature softmax family incur, and how does that cost depend on class count?} Preserving the label ordering within each example does not preserve the ordering of true-label probabilities across examples. It is the latter ordering that determines a LAC cutoff.

We answer this question for a fixed classifier and nonrandomized LAC, whose score is $1-p_y$ \citep{sadinle2018}. The answer is neither an unrestricted overfitting pathology nor a class-count-independent one-parameter bound. The temperature-and-threshold family has VC dimension asymptotic to $\log_2 K$. At an interior finite-rank target and under iid sampling, its worst-selector marginal coverage deficit is, uniformly in $n$ and $K\ge3$,
\[
 \Theta\!\left(\min\left\{1,\sqrt{\frac{\log K}{n}}\right\}\right).
\]
The constants may depend on how far the target is from zero and one. The precise statement, including nominal-coverage rounding and the binary and full-set boundaries, appears in Theorem~\ref{thm:sharp} and Corollary~\ref{cor:nominal}.

\paragraph{Contributions and proof structure.}
First, an exponential-crossing argument bounds the growth function by $m+1+(K-2)m(m-1)$, with all degeneracies and threshold ties allowed. Second, a local realization theorem embeds $K-1$ independently varying coordinates into the true-label probabilities of one fixed classifier at prescribed temperatures. Combining this realization with a classical symmetric-chain decomposition gives a matching leading-order VC lower bound. Third, the same realization transfers independent order statistics into the \emph{actual} selected LAC coverage. A finite-sample Beta-minimum argument then matches the continuum upper rate, including its constant-order saturation. The construction is possible with bounded logits on arbitrarily restricted open temperature intervals, but its numerical margins need not be large.

Finite-family selection bounds already contain square-root logarithmic penalties \citep{yang2024}. General definability results already imply quadratic growth in sample size for a fixed two-parameter formula \citep{aschenbrenner2011}. These facts do not by themselves give a coefficient uniform in the number of softmax classes or realize independent calibration coordinates within one fixed classifier. Our contribution is this explicit structural dependence and its coverage consequence, rather than a new general VC or extreme-value theorem. Section~\ref{sec:related} compares the relevant results in detail.

\paragraph{Scope.}
The lower bound uses an adversarial, calibration-only selector and unmodified LAC sets, which may be empty. It is not a failure theorem for negative-log-likelihood (NLL) tuning. Exact-valid post-selection procedures change the construction and are not contradicted. The theorem concerns worst-case statistical capacity in real arithmetic, not the frequency of failure in applications or the computational cost of a recommended tuning algorithm.

\section{Setup and distinct coverage quantities}\label{sec:setup}
Let $K\ge2$. An input $X$ in a standard Borel space is mapped to finite logits $z(X)\in\R^K$ by a Borel-measurable classifier fixed independently of calibration. We may condition on independent training data throughout. For inverse temperature $\beta=1/T>0$ and label $y\in\{1,\ldots,K\}$, define
\begin{equation}\label{eq:softmax}
 p_\beta(y\mid z)=\frac{e^{\beta z_y}}{\sum_{j=1}^K e^{\beta z_j}},
 \qquad S_\beta(z,y)=1-p_\beta(y\mid z).
\end{equation}
Let $D=(Z_1,\ldots,Z_n)$ with $Z_i=(X_i,Y_i)$ be iid, and let $Z_{n+1}$ be an independent observation with the same law. Fix an integer $1\le k\le n+1$ and write $N=n+1$ and $p=k/N$. A measurable selector chooses $\widehat\beta(D)>0$. For $k\le n$, let $q_D(\beta)$ be the $k$th smallest of the $n$ scores $S_\beta(Z_i)$. The unmodified same-data LAC set is
\begin{equation}\label{eq:set}
 C_D(x)=\{y:S_{\widehat\beta(D)}(z(x),y)\le q_D(\widehat\beta(D))\}.
\end{equation}
For $k=n+1$, the cutoff is $+\infty$ and $C_D(x)$ is the full label set. All claims about Beta order statistics below require $k\le n$.

Write $[x]_+=\max\{x,0\}$. Conditional coverage is $\Cov(D)=\Prob\{Y_{n+1}\in C_D(X_{n+1})\mid D\}$. We distinguish the expected conditional shortfall $\E\pos{p-\Cov(D)}$ from the marginal deficit $\pos{p-\E\Cov(D)}$. Jensen's inequality bounds the latter by the former, but they are not equal. Even a fixed valid score usually has a positive expected conditional shortfall. Neither quantity is automatically the difference of absolute coverage gaps called tuning bias by \citet{zeng2025}.

For a requested nominal coverage $c\in(0,1)$, ordinary conformal indexing sets $k=\lceil Nc\rceil$. Its finite-rank reference level is $p=k/N\in[c,c+1/N)$. A fixed temperature has marginal coverage at least $p$; for continuous true-label scores, it is exactly $p$. Results at $p$ translate to nominal coverage with the explicit rounding correction in Corollary~\ref{cor:nominal}.

We study the indicator class
\begin{equation}\label{eq:class}
 \HH_K=\bigl\{(z,y)\mapsto\1\{p_\beta(y\mid z)\ge t\}:\beta>0,\ t\in[0,1]\bigr\}.
\end{equation}
Its growth function $\Pi_K(m)$ is the maximum number of binary patterns on $m$ points; $d_K=\VC(\HH_K)$ is the largest $m$ for which $\Pi_K(m)=2^m$. A model-specific class is a restriction of this class and can only have smaller complexity. Logarithms without a subscript are natural.

\section{Structural complexity of the temperature family}\label{sec:structure}
\begin{theorem}[Explicit growth and VC bounds]\label{thm:structural}
For every $K\ge2$ and integer $m\ge1$,
\begin{equation}\label{eq:growth}
 \Pi_K(m)\le\min\{2^m,\ m+1+(K-2)m(m-1)\}.
\end{equation}
Define $b(M)=\max\{m\ge1:\binom{m}{\lfloor m/2\rfloor}\le M\}$. Then
\begin{equation}\label{eq:vcbounds}
 b(K-1)\le d_K\le\log_2 K+2\log_2\log_2 K+O(1)
 \quad(K\to\infty).
\end{equation}
In particular,
\begin{equation}\label{eq:vcasymptotic}
 d_K=(1+o(1))\log_2 K,
 \qquad b(K-1)=\log_2 K+\tfrac12\log_2\log_2 K+O(1).
\end{equation}
The weaker lower bound $\lfloor\log_2(K-1)\rfloor$ also follows directly. Moreover, $d_2=1$ and $d_3=2$. The lower bounds remain valid when inverse temperatures are restricted to any prescribed nonempty open interval $I\subset(0,\infty)$ and logits to $[-B,0]^K$ for any $B>0$.
\end{theorem}
The upper bound is distribution-free. The restricted lower bound is existential and has no margin lower bound uniform in $K$, $I$, or $B$.

\subsection{Bounding cross-example score crossings}
For a sampled point $(z_i,y_i)$, put $c_{ij}=z_{ij}-z_{iy_i}$ and
\begin{equation}\label{eq:denom}
 A_i(\beta)=p_\beta(y_i\mid z_i)^{-1}=1+\sum_{j\ne y_i}e^{\beta c_{ij}}.
\end{equation}
A nonzero linear combination of $r$ exponentials with distinct real exponents has at most $r-1$ distinct real zeros. Thus two nonidentical curves in~\eqref{eq:denom} have at most $2K-4$ positive intersections: the difference has at most $2(K-1)$ terms, and the common zero at $\beta=0$ uses one of the available zeros. Adding the graphs one at a time gives at most $m+1+(2K-4)\binom m2$ open regions. Threshold ties can be moved into an open region without changing the pattern. Appendix~\ref{app:growth} proves this carefully, including degenerate intersections.

\subsection{One fixed classifier, multiple score coordinates}
\begin{theorem}[Local softmax realization]\label{thm:realization}
Let $M=K-1$. Fix any distinct positive numbers $\beta_1,\ldots,\beta_M$ and $B>0$. There are constants $a_1,\ldots,a_M>1$, a common $\delta>0$, and a smooth map $z$ on a neighborhood of $[0,1]^M$ into $[-B,0]^K$ such that $z_1(u)=0$, $-B<z_j(u)<0$ for $j>1$, and
\begin{equation}\label{eq:realization}
 \frac{1}{p_{\beta_r}(1\mid z(u))}=a_r+\delta u_r
 \quad\text{for all }u\in[0,1]^M,\ r=1,\ldots,M.
\end{equation}
The map and constants are fixed before calibration, and may be chosen independently of $n$ and $k$.
\end{theorem}
To see the mechanism, choose distinct $c_j^0\in(-B,0)$ and let $F_r(c)=\sum_j e^{\beta_r c_j}$. The Jacobian $J_{rj}=\beta_r e^{\beta_r c_j^0}$ is nonsingular: a linear dependence would create an $M$-term exponential sum with $M$ distinct zeros. The inverse-function theorem supplies a local inverse. A sufficiently small translated cube fits inside its open image, giving~\eqref{eq:realization}. This is local interpolation by a fixed logit map, not post-calibration modification of the classifier.

The lower bound in Theorem~\ref{thm:structural} uses these coordinates to produce a family of score orderings. The classical symmetric-chain partition of the subsets of $m$ points \citep{debruijn1951} has $\binom{m}{\lfloor m/2\rfloor}$ chains. Each chain extends to the prefixes of a permutation. One realized coordinate supplies each permutation, and a cutoff chooses its prefixes. Every subset is therefore obtained. The elementary chain construction and the full local-inverse argument are in Appendix~\ref{app:realization}.

\section{Coverage upper bound for arbitrary reuse}\label{sec:upper}
\begin{theorem}[Expected conditional shortfall]\label{thm:upper}
Under the iid and independence assumptions in Section~\ref{sec:setup}, every measurable temperature selector satisfies, for a universal constant $C$,
\begin{equation}\label{eq:upper}
 \E\pos{p-\Cov(D)}\le\min\{1,C\sqrt{d_K/n}\}
 \le\min\{1,C'\sqrt{\log K/n}\}.
\end{equation}
The constants do not depend on the fixed model, input distribution, class count, selector, or $k$. Ties are allowed. For $k=n+1$ the left side is zero.
\end{theorem}
For $k\le n$, the selected rule $h_D\in\HH_K$ includes at least $k$ calibration observations, so $P_nh_D\ge k/n\ge p$. Consequently,
\begin{equation}\label{eq:empirical}
 \pos{p-Ph_D}\le\sup_{h\in\HH_K}|P_nh-Ph|.
\end{equation}
A standard expected VC empirical-process bound gives $C\sqrt{d_K/n}$. This background bound is not a new result; it is also the content of \citet[Lemma D.6]{zeng2025}. Appendix~\ref{app:empirical} states the inequality precisely, addresses measurability, and gives an entropy proof sketch. It also provides an explicit but looser finite-trace bound carrying an additional $\log n$ factor. We do not treat an unspecified $C$ in~\eqref{eq:upper} as an operational calibration correction.

\section{An exact coverage law and a uniform matching lower bound}\label{sec:lower}
\begin{theorem}[Coverage law inside a fixed softmax model]\label{thm:law}
For every $K\ge2$, the fixed classifier in Theorem~\ref{thm:realization} and the fixed distribution $U\sim\operatorname{Uniform}([0,1]^{K-1})$, $Y=1$, admit a measurable temperature selector with the following property. For $1\le k\le n$, put $M=K-1$. Its conditional coverage has the law
\begin{equation}\label{eq:minlaw}
 \Cov(D)=\min_{1\le r\le M} B_r,
 \qquad B_r\overset{\mathrm{iid}}\sim\Beta(k,n+1-k).
\end{equation}
Thus its marginal coverage is
\begin{equation}\label{eq:integral}
 \mu_{n,k,M}=\int_0^1\bigl(1-F_{\Beta(k,n+1-k)}(t)\bigr)^M\,dt.
\end{equation}
Every fixed candidate temperature has marginal coverage $p=k/(n+1)$. For $k=n+1$, all candidates and the selected rule instead have coverage one; no Beta distribution with a zero shape parameter is used.
\end{theorem}
At temperature $\beta_r$, the true-label score is a strictly increasing function of $U_r$. Its empirical cutoff corresponds to the $k$th order statistic $B_r$ of that coordinate. The selector chooses the smallest \emph{normalized} cutoff $B_r$, with a fixed tie rule. Normalization is available from the known constants in~\eqref{eq:realization}: $B_r=((1-q_D(\beta_r))^{-1}-a_r)/\delta$. This is a legal calibration-only rule but is intentionally adversarial; it does not minimize NLL or necessarily minimize set size. Independence follows from the product distribution across coordinates, not from treating repeated measurements on the same observation as independent without justification.

\begin{lemma}[Uniform loss of the minimum]\label{lem:uniform}
Let $0<\eta\le1/2$, $N=n+1\ge16/\eta^2$, $p=k/N\in[\eta,1-\eta]$, and $M\ge2$. There is $L_\eta>0$, depending only on $\eta$, such that
\begin{equation}\label{eq:uniformmin}
 p-\mu_{n,k,M}\ge L_\eta\min\{1,\sqrt{\log M/N}\}.
\end{equation}
This is simultaneous in $n,k,M$ in the stated regime, including arbitrarily large $M$.
\end{lemma}
The lemma is a consequence of established Beta tail behavior, not a new general probability principle; compare \citet[Theorem 8]{zhang2020}. Appendix~\ref{app:beta} supplies a self-contained derivation with explicit, conservative constants. The useful device is to write the loss as a nonnegative difference $\E[B_1-\min_r B_r]$, avoiding cancellation from a two-sided tail integral.

\begin{definition}[Worst-selector deficit]\label{def:risk}
Let $\risk_{n,K,k}$ be the supremum of $\pos{p-\E\Cov(D)}$ over Borel fixed classifiers, iid input--label laws, and measurable positive-temperature selectors using~\eqref{eq:set}. Define $\risk^{I,B}_{n,K,k}$ in the same way with $\widehat\beta\in I$ and logits in $[-B,0]^K$, where $I\subset(0,\infty)$ is nonempty and open and $B>0$.
\end{definition}
\begin{theorem}[Uniform sharp rate at the finite-rank level]\label{thm:sharp}
There are a universal $A>0$ and, for every $0<\eta\le1/2$, a constant $a_\eta>0$, such that for every $K\ge3$, $n+1\ge16/\eta^2$, and $p=k/(n+1)\in[\eta,1-\eta]$,
\begin{equation}\label{eq:sharp}
 a_\eta\min\{1,\sqrt{\log K/n}\}
 \le\risk^{I,B}_{n,K,k}
 \le\risk_{n,K,k}
 \le A\min\{1,\sqrt{\log K/n}\}.
\end{equation}
This holds for every $I,B$ in Definition~\ref{def:risk}, with constants independent of $I,B$. For $K=2$, $\risk_{n,2,k}=0$ for every $n,k$. For $k=n+1$, the deficit is zero for every $K$.
\end{theorem}
The upper bound is Theorem~\ref{thm:upper}. The lower bound uses Theorems~\ref{thm:realization} and~\ref{thm:law}, Lemma~\ref{lem:uniform}, $M=K-1$, $\log(K-1)\ge\tfrac12\log K$ for $K\ge3$, and $n+1\le2n$. One may take $a_\eta=L_\eta/2$. Binary invariance is proved in Appendix~\ref{app:small}. This is a characterization over the worst selector, not a minimax impossibility over all possible valid reuse procedures.

\begin{corollary}[Nominal coverage and rounding]\label{cor:nominal}
For $k=\lceil(n+1)c\rceil\le n$, suppose $p=k/(n+1)\in[\eta,1-\eta]$ and $n+1\ge16/\eta^2$. The constructed rule satisfies
\begin{equation}\label{eq:rounding}
 \pos{c-\E\Cov(D)}\ge
 \pos{L_\eta\min\{1,\sqrt{\log(K-1)/(n+1)}\}-\frac1{n+1}}.
\end{equation}
For each fixed interior range $c\in[\eta_0,1-\eta_0]$, the worst nominal deficit has the same capped rate for all $K\ge3$ and all sufficiently large $n$, with the threshold and lower constant depending only on $\eta_0$.
\end{corollary}
The $1/(n+1)$ subtraction must not be silently omitted. For the last assertion use $\eta=\eta_0/2$ and note that the smallest rate over $K\ge3$ is at least $\sqrt{\log2/(n+1)}$; eventually rounding is at most half the lower bound. The resulting constants are not practical sample-size recommendations.

\paragraph{A complementary limit.}
For fixed $M$, $k=\lceil(n+1)c\rceil$, and $c\in(0,1)$,
\begin{equation}\label{eq:iterated}
 \sqrt n\,(p-\mu_{n,k,M})\longrightarrow
 \sqrt{c(1-c)}\,\E\max_{r\le M}G_r,
 \quad G_r\overset{\mathrm{iid}}\sim N(0,1).
\end{equation}
Appendix~\ref{app:clt} proves expectation convergence. Subsequently increasing $M$ gives the usual $\sqrt{2\log M}$ factor. This iterated limit is distinct from the simultaneous finite-sample statement proved above.

\section{A three-class illustration and finite certificates}\label{sec:evidence}
Theoretical statements above do not depend on simulations. Computation has two roles: finite exact certificates that test delicate orderings, and a small implementation illustration of the coverage construction. No training or real-data claim is needed for these theoretical results.

\subsection{An explicit three-class model}
For $U_1,U_2$ independent uniform variables, define
\begin{equation}\label{eq:explicit}
 s=1+0.01U_1,\quad v=\tfrac58+0.01U_2,\quad
 w_\pm=\frac{s\pm\sqrt{2v-s^2}}2,\quad z=(0,\log w_-,\log w_+).
\end{equation}
The square root is real and $0<w_-<w_+<1$ throughout the unit square. At inverse temperatures $1$ and $2$,
\begin{equation}
 p_1(1\mid z)^{-1}=2+0.01U_1,
 \qquad p_2(1\mid z)^{-1}=\tfrac{13}{8}+0.01U_2.
\end{equation}
For $n=99$, $k=90$, and $M=2$, exact rational polynomial integration gives
\begin{equation}\label{eq:exactnumber}
 \mu_{99,90,2}=\frac{5214610247405709}{5903643345897130}
 =0.883286801366434\ldots.
\end{equation}
Numerical quadrature agrees within $2.23\times10^{-16}$. The fraction, not the quadrature error estimate, is the exact certificate.

\begin{table}[t]
\centering\small
\caption{Three-class illustration, 10,000 iid calibration experiments, nominal coverage 90\%. Entries are percentages. Conditional fresh-test coverage is integrated analytically for each fitted set. Intervals quantify Monte Carlo uncertainty only.}\label{tab:coverage}
\begin{tabular}{lrrrr}
\toprule
Procedure & Observations & Expected & Measured & 95\% MC interval\\
\midrule
Same-data adversarial reuse & 99 & 88.3287 & 88.2848 & [88.2313, 88.3383]\\
Fixed temperature & 99 & 90.0000 & 89.9480 & [89.8896, 90.0063]\\
Independent tuning/calibration & 198 & 90.0000 & 90.0474 & [89.9889, 90.1060]\\
Same-total-data split & 99 & 90.0000 & 89.9411 & [89.8580, 90.0242]\\
Analytic NLL grid control & 99 & 90.0000 & 89.9965 & [89.9377, 90.0554]\\
\bottomrule
\end{tabular}
\end{table}

Table~\ref{tab:coverage} uses 10,000 independent calibration samples with fixed seeds. The same-total-data control uses 50 tuning and 49 calibration observations, with rank $45$. The independent control uses 99 observations at each stage. The implementation computes softmax scores and conformal quantiles; conditional fresh-test coverage is then integrated analytically. The NLL control uses its analytically determined choice $\beta=2$, rather than running a numerical optimizer. The supplementary material gives raw outputs, uncertainty calculations, and reproduction instructions.

\paragraph{Negative controls delimit the example.}
All true labels are the unique argmax. Increasing $\beta$ improves their likelihood for every observation, so the NLL grid selector always chooses $\beta=2$ and does not under-cover in this construction. Adding the argmax label to every prediction set likewise eliminates this example's failures. The theorem concerns ordinary LAC, which permits empty sets, not that modified rule.

\subsection{Exact ordering certificates}
Finite certificates instantiate the symmetric-chain construction using rational linear algebra and an analytic Taylor remainder. They establish eight patterns on three points with $K=4$ and sixteen patterns on four points with $K=7$, hence $d_4\ge3$ and $d_7\ge4$. These are lower bounds, not exact dimensions at those class counts. Appendix~\ref{app:certificates} describes the certificates. An earlier floating-point candidate failed exact verification; it is excluded from every claim and retained in the supplementary negative-result record.

\section{Related work}\label{sec:related}
\paragraph{Temperature and conformal prediction.}
\citet{dabah2025} study the effect of temperature on conformal efficiency and coverage trade-offs, particularly for adaptive classification scores. \citet{mattei2026} develop elementary temperature-scaling properties, including within-example ordering and likelihood structure. Neither within-example label-order preservation nor the existence of a temperature trade-off is claimed here. Our object is the set of true-label score orderings \emph{across examples}, together with an adaptive LAC cutoff.

\paragraph{Complexity-based reuse bounds.}
\citet{zeng2025} analyze tuning bias through an empirical-process discrepancy and establish binary order preservation. Their Lemma D.6 supplies the standard VC transfer. However, the parameter-count argument in arXiv v2, Lemmas D.3--D.5, does not follow from continuity and boundedness alone. Rather than using those assumptions as a capacity bound, we calculate the complexity of the specific temperature family. The growth of $d_K$ in Theorem~\ref{thm:structural} rules out a class-count-independent VC bound for this family. This does not refute their empirical observations about particular tuning rules.

\paragraph{Reuse, model selection, and validity.}
\citet[Theorem 1]{yang2024} give finite-family coverage costs and also study additional splitting. \citet[Theorems 3.1 and 3.3]{liang2026} obtain exact-valid post-selection conformal procedures by restoring suitable symmetry. \citet{wang2026} use localized model selection, and \citet{hegazy2025} study stability-based selection guarantees. These results address procedures or restrictions absent from our adversarial unmodified rule. They remain alternatives, not targets of an impossibility claim. We do not establish a new algorithmic advantage over them.

\paragraph{Complexity of exponential families.}
\citet{karpinski1997} bound VC complexity in general Pfaffian and sigmoidal settings using geometry of parameter-space sign conditions. \citet[Theorem 1.1]{aschenbrenner2011} bound VC density by the number of parameters in weakly o-minimal theories. For a fixed $K$ in the real exponential structure, the latter implies an $O_K(m^2)$ growth bound. Quadratic growth for a fixed two-parameter definable class is therefore not new. Our explicit coefficient in~\eqref{eq:growth}, class-count VC asymptotics, fixed-model realization, and coverage transfer retain the class-count information that the fixed-formula result leaves implicit. The proof uses classical exponential zero counting, local inversion, and symmetric chains; we do not claim those tools as discoveries.

\paragraph{Probability and empirical-process background.}
Expected VC inequalities, obtained from entropy and symmetrization, are standard \citep{haussler1995}. Beta lower tails and logarithmic extreme-value effects are also established \citep{zhang2020}. The probability lemma here is included for a self-contained finite-sample transfer. The softmax-specific conclusion is that this behavior is available inside a single, bounded-logit, globally temperature-scaled classifier and is matched by a continuum upper bound. The independent-coordinate realization is the link between the constrained score family and the probability calculation.

\section{Limitations and interpretation}
The asymptotically leading VC coefficient and the finite-rank worst-case rate are precise, but several limits matter. The exact VC dimension is not determined for every fixed $K$; the second-order upper and lower terms in~\eqref{eq:vcbounds} do not match. The realization radius can be extremely small. Compact temperature ranges and bounded logits do not control finite precision, quantization, or probability margins. We provide no uniform bit-complexity guarantee.

The lower selector uses a normalized empirical cutoff constructed to expose the effect. It is neither NLL tuning nor a demonstrated practical set-size optimizer. No real-data prevalence, efficiency gain, set-size benefit, or benchmark reversal has been established. The generic upper bound can be conservative, and its unspecified universal constant is not an operational correction. An algorithm-specific guarantee would require additional analysis, such as usable stability or likelihood-curvature conditions, and cannot be read off from the worst-selector theorem.

Our iid setting is narrower than general exchangeability. The model must be fixed independently of calibration; joint retraining, covariate shift, vector scaling, feature-dependent temperature, randomized scores, and forced-argmax LAC are outside the claims. 

\section{Conclusion}
For ordinary temperature-scaled LAC, the relevant capacity is controlled but class-count dependent. A global scalar temperature and an adaptive cutoff can express enough cross-example orderings to give VC dimension asymptotic to $\log_2 K$. A fixed-model interpolation converts that capacity into an exact coverage law and a uniform, capped square-root-logarithmic worst-selector deficit. This delineates the strongest distribution-free rate available for the unmodified procedure while leaving ordinary likelihood tuning and exact-valid reuse methods as distinct questions.

\section*{Reproducibility statement}
The supplementary code reproduces the numerical illustration, exact finite certificates, and boundary checks without model training or inference services. Complete proofs are included below; simulations are not used to justify universal statements. The supplement distinguishes rational certificates, deterministic numerical integration, and Monte Carlo results, and preserves failed witnesses and negative controls.

\section*{Use of generative AI}
Generative AI tools assisted mathematical exploration, code development, proof review, and manuscript drafting. No independent external mathematical validation is claimed.

\bibliographystyle{tmlr}
\bibliography{references}

\clearpage
\appendix
\section{Zero counts, arrangements, and the VC upper bound}\label{app:growth}
\begin{lemma}[Real exponential zero count]\label{lem:zeros}
If $a_1<\cdots<a_r$ and not all $b_j$ vanish, then $f(x)=\sum_{j=1}^r b_je^{a_jx}$ has at most $r-1$ distinct real zeros.
\end{lemma}
\begin{proof}
Remove zero coefficients. The assertion is immediate for one term. Divide by $e^{a_1x}$, which does not change zeros. The derivative of the resulting function has at most $r-1$ nonzero exponential terms with distinct exponents, hence at most $r-2$ zeros by induction. If the original function had $r$ distinct zeros, Rolle's theorem between consecutive zeros would yield at least $r-1$ derivative zeros, a contradiction. Applying this argument to any finite selection of zeros also rules out infinitely many zeros. This proof counts distinct zeros; no unproved multiplicity claim is used.
\end{proof}

\begin{proof}[Proof of the growth bound]
For each pair $i,j$, combine equal exponents in $A_i-A_j$. Either the curves are identical, or the difference has at most $2(K-1)$ terms and therefore at most $2K-3$ distinct real zeros. The zero at $0$ always exists since $A_i(0)=A_j(0)=K$. At most $2K-4$ zeros remain on $(0,\infty)$. Negative zeros can only reduce this number.

Merge identical curves; they cannot independently label sampled points and cannot increase the number of patterns. In the strip $(0,\infty)\times\R$, insert the remaining graphs successively. If a new graph has $q$ distinct intersections with the preceding union, its complement of the intersection points has $q+1$ connected arcs. Adding each arc increases the number of connected faces by at most one. This follows by restricting to a large rectangle containing any finite configuration under consideration, or by the elementary planar separation property of a simple arc. Therefore the number of faces is at most one plus the number of graphs plus the total pairwise intersection count. A tangency may add no face; simultaneous crossings reduce the number of distinct split points. Counting them pairwise is an upper bound in both cases.

For $t>0$, a pattern is $\{i:A_i(\beta)\le u\}$ with $u=1/t$. If $u$ equals one or more $A_i(\beta)$, increase $u$ by less than the minimum positive gap to an excluded value. All included equalities remain included and no excluded point enters. If nothing is excluded, take any larger finite $u$. All inequalities are now strict; a sufficiently small change in $(\beta,u)$ retains them, so the pattern occurs in an open face. The empty pattern also occurs strictly below all curves. The endpoint $t=0$ produces the full pattern, already counted; $t=1$ is empty because all logits are finite and $K\ge2$. Using the full vertical strip rather than $u\ge1$ only enlarges the count. Each open face has a fixed pattern, proving~\eqref{eq:growth}.
\end{proof}

\begin{proof}[VC upper asymptotics]
For a shattered set of $d\ge2$ points, $2^d\le d+1+(K-2)d(d-1)\le Kd^2$. The last inequality holds for $K\ge2,d\ge2$. Writing $L=\log_2K\ge1$, the function $x-2\log_2x$ is increasing for $x\ge4$ and its value at $4(L+1)$ exceeds $L$. Hence $d<4(L+1)$, in particular $d=O(\log K)$. Taking base-two logarithms gives $d\le\log_2K+2\log_2 d$, and substitution of $d=O(\log K)$ proves the claimed upper expansion. The finitely many small $K$ are absorbed into the constant when using $d_K\le C_0\log K$ in a uniform inequality.
\end{proof}

\section{Realization and the refined combinatorial lower bound}\label{app:realization}
\begin{proof}[Proof of Theorem~\ref{thm:realization}]
Choose distinct $c_j^0\in(-B,0)$, $j=1,\ldots,M$, and consider the analytic map
\[
 F:\R^M\to\R^M,\qquad F_r(c)=\sum_{j=1}^M e^{\beta_r c_j}.
\]
If $Jv=0$ at $c^0$, division of row $r$ by $\beta_r>0$ gives $\sum_jv_je^{\beta_rc_j^0}=0$. For nonzero $v$, Lemma~\ref{lem:zeros} allows at most $M-1$ distinct zeros as a function of $\beta$, contradicting the $M$ prescribed distinct values. Hence $J$ is invertible.

The inverse-function theorem supplies an open neighborhood $V\subset(-B,0)^M$ of $c^0$ on which $F$ is an analytic diffeomorphism onto an open neighborhood $W$ of $F(c^0)$. Choose $\delta>0$ so that a neighborhood of $F(c^0)+\delta[0,1]^M$ is contained in $W$. This is possible because the cube has finite dimension and $W$ is open: for example, fit a cube inside an open Euclidean ball about $F(c^0)$. Set
\[
 c(u)=F^{-1}(F(c^0)+\delta u),\quad
 z(u)=(0,c_1(u),\ldots,c_M(u)),\quad a_r=1+F_r(c^0).
\]
These definitions prove all identities, bounds, and smoothness claims. The common $\delta$ may depend on all prescribed temperatures and on $B,K,c^0$. No lower bound on it independent of these choices is asserted. All choices precede the sample. A prescribed open interval contains any finite number of distinct temperatures, proving the restricted-interval assertion.
\end{proof}

\begin{lemma}[Symmetric chains; \citealp{debruijn1951}]\label{lem:chains}
The subsets of an $m$-element set can be partitioned into $\binom m{\lfloor m/2\rfloor}$ saturated chains, each with endpoint cardinalities summing to $m$.
\end{lemma}
\begin{proof}
Start with the single chain $\{\varnothing\}$ for $m=0$. Suppose a chain for $m$ is $A_r\subset A_{r+1}\subset\cdots\subset A_{m-r}$, where subscripts denote cardinalities. On adjoining a new element $e$, replace it by
\[
 A_r\subset\cdots\subset A_{m-r}\subset A_{m-r}\cup\{e\}
\]
and, when nonempty, the chain
\[
 A_r\cup\{e\}\subset\cdots\subset A_{m-r-1}\cup\{e\}.
\]
These partition the original chain and its copy containing $e$, are saturated, and have endpoint sizes summing to $m+1$. Induction gives the partition. Each symmetric chain contains exactly one set of size $\lfloor m/2\rfloor$; counting that layer gives the number of chains.
\end{proof}

\begin{proof}[Proof of the structural lower bounds]
Fix $m$ with $q=\binom m{\lfloor m/2\rfloor}\le M=K-1$. Extend each of the $q$ chains to the prefixes of a permutation of $\{1,\ldots,m\}$. For point $i$ and coordinate $r\le q$, let $u_r^{(i)}$ be its rank in permutation $r$, scaled into $[0,1]$; for $m=1$ choose any value. Set unused coordinates to zero. Apply Theorem~\ref{thm:realization} to these fixed points. At $\beta_r$, reciprocal probabilities are $a_r+\delta u_r^{(i)}$, so varying the cutoff obtains all prefixes of permutation $r$. Every subset belongs to one original chain, and hence is a prefix of at least one extended permutation. Empty and full subsets are obtainable by extreme cutoffs. Thus all $2^m$ patterns occur and $d_K\ge b(M)$.

An alternative weaker construction assigns one coordinate to each binary pattern and sets its values to zero on included points and one on excluded points, obtaining $d_K\ge\lfloor\log_2M\rfloor$. For the refinement, the central-binomial estimate
$\binom m{\lfloor m/2\rfloor}=\Theta(2^m/\sqrt m)$ implies
$b(M)=\log_2M+\tfrac12\log_2\log_2M+O(1)$; this follows by taking logarithms in the defining inequality and its failure at $b(M)+1$. Combining with the upper bound proves the leading coefficient. The realization uses bounded logits in any prescribed interval of temperatures, so the same lower arguments apply under those restrictions.
\end{proof}

\section{Small-class boundaries}\label{app:small}
\begin{proposition}\label{prop:small}
For $K=2$, $d_2=1$, and same-data temperature selection does not change the LAC set at all. For $K=3$, $d_3=2$.
\end{proposition}
\begin{proof}
For $K=2$, let $a=z_{\mathrm{other}}-z_y$. The score $S_\beta=e^{\beta a}/(1+e^{\beta a})$ is a strictly increasing function of $a$ for every $\beta>0$. The $k$th calibration score is this same transform of the $k$th calibration gap. Test inclusion is therefore exactly $a_{\rm test}\le a_{(k)}$, independent of the selected temperature, including ties. The marginal rank argument gives coverage at least $k/(n+1)$, proving zero marginal deficit. The class consists of thresholds in $a$, so its VC dimension is one. Conditional coverage remains random; expected conditional shortfall need not vanish.

For $K=3$, write the two competitor gaps as $a,b$ and let $u=(a+b)/2$, $v=|a-b|/2$. The denominator is $1+2e^{\beta u}\cosh(\beta v)$. Equality of two nonidentical denominators is equivalent to
\[
 u_1-u_2+\frac{\log\cosh(\beta v_1)-\log\cosh(\beta v_2)}\beta=0.
\]
The derivative of the fraction has the sign of $g(\beta v_1)-g(\beta v_2)$, where
$g(x)=x\tanh x-\log\cosh x$ and $g'(x)=x\operatorname{sech}^2x$. If $v_1\ne v_2$, this sign is fixed and nonzero for $\beta>0$, so there is at most one positive crossing. If $v_1=v_2$, the curves are either identical or never cross. The arrangement bound is now $m+1+\binom m2$, which is seven at $m=3$, proving $d_3\le2$. The lower bound follows from Lemma~\ref{lem:chains} and the two-coordinate realization. Independently, gaps $(-1,-10)$ and $(-4,-4)$ have opposite denominator order at $e^\beta=11/10$ and $2$, with exact differences $-1851743090/25937424601$ and $385/1024$. They realize both singleton patterns, with empty and full patterns obtained by the cutoff.
\end{proof}

\section{Measurability and the empirical-process bound}\label{app:empirical}
The upper bound uses a pointwise measurable indicator class. This can be established directly rather than assumed from its parameterization. Work with $h_{\beta,u}=\1\{A_\beta\le u\}$ and adjoin the two constant functions. For any $(\beta,u)$ choose rational $\beta_j$ within $j^{-2}$ of $\beta$, and rational $u_j\in(u+j^{-1},u+2j^{-1})$. For each fixed finite logit vector, $A_\beta$ has a locally bounded derivative. Its change under $\beta_j$ is $O(j^{-2})$, smaller than the positive threshold slack. Equality at $A_\beta=u$ is preserved eventually, as are strict inequalities. Therefore $h_{\beta_j,u_j}$ converges pointwise to $h_{\beta,u}$. Dominated convergence gives convergence of population expectations and of empirical averages. Consequently the supremum in~\eqref{eq:empirical} equals a supremum over a countable subclass and is measurable. On a restricted open temperature interval the rationals may also be chosen inside that interval.

For a pointwise measurable $\{0,1\}$-valued class of VC dimension $d\ge1$, the background inequality is
\begin{equation}\label{eq:vcbackground}
 \E\sup_{h\in\HH}|(P_n-P)h|\le C\sqrt{d/n}.
\end{equation}
For clarity, its standard route is symmetrization followed by an entropy integral. The VC packing bound gives, uniformly in probability measures $Q$,
$N(\epsilon,\HH,L_2(Q))\le c_1(c_2/\epsilon)^{c_3d}$ for universal constants and $0<\epsilon<1$ \citep{haussler1995}. Conditioning on the observations and applying the Rademacher entropy integral yields at most a constant times
\[
 \frac1{\sqrt n}\int_0^1\sqrt{d\log(c_2/\epsilon)}\,d\epsilon
 =O(\sqrt{d/n}).
\]
Symmetrization changes only the universal constant. This standard expected bound is also stated in \citet[Lemma D.6]{zeng2025}. A union bound on all empirical patterns would instead introduce an unnecessary $\log n$ factor.

For readers needing a completely explicit, looser expression, define
$B_{K,n}=\min\{2^n,n+1+(K-2)n(n-1)\}$. Conditional on the sample there are at most $B_{K,n}$ trace vectors of Euclidean norm at most $\sqrt n$. The finite-class Rademacher maximum bound and symmetrization give
\begin{equation}\label{eq:explicitupper}
 \E\sup_h|(P_n-P)h|\le 2\sqrt{\frac{2\log(2B_{K,n})}{n}}.
\end{equation}
Changing one observation changes the supremum by at most $1/n$. Bounded differences therefore adds $\sqrt{\log(1/\varepsilon)/(2n)}$ to~\eqref{eq:explicitupper} for a probability-$1-\varepsilon$ bound. Values may always be capped at one. The sharper asymptotic rate and this explicit but looser constant bound are distinct statements.

\section{Coverage law and uniform Beta-minimum loss}\label{app:beta}
\begin{proof}[Proof of Theorem~\ref{thm:law}]
Use the fixed model from Theorem~\ref{thm:realization} and iid vectors $U_i$ with independent uniform coordinates. At $\beta_r$, $S_{\beta_r}=1-(a_r+\delta U_r)^{-1}$ is strictly increasing in $U_r$. Thus its $k$th calibration cutoff maps back exactly to $B_r=U_{(k),r}$. The arrays $\{U_{ir}:i\le n\}$ are independent across $r$, so the $B_r$ are independent $\Beta(k,n+1-k)$ variables. Conditional on $D$, the selected coordinate is fixed and its independent test value is uniform. Its acceptance probability is the selected $B_r$. Choosing an index of the minimum, with deterministic tie breaking, gives~\eqref{eq:minlaw}. The tail integral for a nonnegative bounded random variable gives~\eqref{eq:integral}. All selections and order statistics are Borel measurable. At $k=n+1$ the cutoff is instead $+\infty$, giving full sets.
\end{proof}

\begin{proof}[Proof of Lemma~\ref{lem:uniform}]
We provide deliberately conservative constants. Put
\begin{equation}\label{eq:constants}
 A_\eta=e^{-1/6}\sqrt{\frac{\eta(1-\eta)}{2\pi}},\quad
 B_\eta=\frac4\eta,\quad
 r_\eta=A_\eta e^{-B_\eta},\quad q_\eta=A_\eta e^{-2B_\eta},
\end{equation}
and set
\begin{equation}\label{eq:L}
 L_\eta=\frac\eta4\,r_\eta\bigl(1-e^{-q_\eta/\sqrt2}\bigr)>0.
\end{equation}
For $B\sim\Beta(k,N-k)$ with $p=k/N\in[\eta,1-\eta]$, its density $f$ obeys
\begin{equation}\label{eq:density}
 f(x)\ge A_\eta\sqrt N\exp\{-B_\eta N(x-p)^2\},
 \qquad x\in[\eta/2,1-\eta/2].
\end{equation}
Here are the constants explicitly. The standard Stirling bounds
$\Gamma(s)=\sqrt{2\pi}s^{s-1/2}e^{-s}e^{\rho_s}$, $0<\rho_s<1/(12s)$ for positive integers $s$, give
\[
 f(x)\ge e^{-1/6}\sqrt{\frac N{2\pi}}
 \frac{\sqrt{p(1-p)}}{x(1-x)}e^{-ND(p\Vert x)},
\]
where $D(p\Vert x)=p\log(p/x)+(1-p)\log((1-p)/(1-x))$. The correction exponent is at least $-1/(12k)-1/(12(N-k))\ge-1/6$. The inequality $\log t\le t-1$ implies
$D(p\Vert x)\le(p-x)^2/[x(1-x)]$. On the displayed interval $x(1-x)\ge\eta/4$. Also $p(1-p)\ge\eta(1-\eta)$ and $1/[x(1-x)]\ge1$. These establish~\eqref{eq:density}.

Let $h=N^{-1/2}$ and
\[
 t=\frac\eta4\min\{1,\sqrt{\log M/N}\}.
\]
The sample-size assumption ensures $h\le\eta/4$, and $t\le\eta/4$. The intervals $[p,p+h]$ and $[p-t-h,p-t]$ lie inside $[\eta/2,1-\eta/2]$. Integrating~\eqref{eq:density} over the first gives
\[
 \Prob(B\ge p)\ge A_\eta e^{-B_\eta}=r_\eta.
\]
Over the second, using $(t+h)^2\le2t^2+2h^2$, it gives
\[
 \Prob(B\le p-t)\ge A_\eta e^{-2B_\eta}e^{-2B_\eta Nt^2}
 \ge q_\eta M^{-1/2}.
\]
Indeed $Nt^2\le(\eta^2/16)\log M$, so $2B_\eta Nt^2\le(\eta/2)\log M\le\tfrac12\log M$. For $M\ge2$, $(M-1)/\sqrt M\ge1/\sqrt2$, hence
\[
 \Prob\{\min_{2\le r\le M}B_r\le p-t\}
 \ge1-\exp\{-(M-1)q_\eta M^{-1/2}\}
 \ge1-e^{-q_\eta/\sqrt2}.
\]
The two events involving $B_1$ and the remaining coordinates are independent. Finally,
\[
 p-\E\min_{r\le M}B_r
 =\E\bigl[B_1-\min_{r\le M}B_r\bigr]
 =\E\pos{B_1-\min_{2\le r\le M}B_r}.
\]
On their intersection the last integrand is at least $t$. Multiplication of the probabilities proves~\eqref{eq:uniformmin} with~\eqref{eq:L}. The capped definition of $t$ is why no restriction such as $\log M=o(n)$ is needed. This argument proves existence of positive constants, not a useful numerical calibration penalty.
\end{proof}

\section{Expectation convergence for the iterated limit}\label{app:clt}
For fixed $M$ and $k=\lceil(n+1)c\rceil$, a Beta variable can be represented as $G_k/(G_k+G_{N-k}')$, with independent unit-scale gamma variables having the indicated integer shapes. They are sums of independent unit exponentials. The ordinary joint central limit theorem and the delta method give
\[
 \sqrt n(B_r-p)\ \Rightarrow\ \sqrt{c(1-c)}G_r
\]
jointly for the fixed number of independent coordinates. The continuous mapping theorem applies to their minimum. To pass to expectations, use
\[
 \E\left[\max_{r\le M}|\sqrt n(B_r-p)|^2\right]
 \le\sum_{r=1}^M n\operatorname{Var}(B_r)
 =\frac{Mn\,p(1-p)}{N+1}\le\frac M4.
\]
This uniform second-moment bound gives uniform integrability, proving~\eqref{eq:iterated}. For completeness, the usual normal-maximum asymptotics follow by the exponential-moment upper bound $\E\max_rG_r\le\sqrt{2\log M}$ and the normal-tail lower bound: for any $\epsilon>0$, $M\Prob\{G>(1-\epsilon)\sqrt{2\log M}\}\to\infty$. Thus the maximum exceeds this level with probability tending to one. Its negative part has expectation tending to zero, for example by integrating $\Prob(\max_rG_r<-t)=\Phi(-t)^M$. Dividing by $\sqrt{2\log M}$ and letting $\epsilon\downarrow0$ gives the limit. None of these fixed-$M$ expectation arguments are substituted for the uniform finite-sample proof.

\section{Exact finite certificate methods}\label{app:certificates}
\subsection{Rational integration of the primary coverage law}
For integer $k\le n$, put $b=n+1-k$. The Beta density expands as
\[
 f(t)=\frac{n!}{(k-1)!(b-1)!}\sum_{j=0}^{b-1}(-1)^j\binom{b-1}{j}t^{k+j-1}.
\]
Integrating term by term gives an exact rational polynomial for $F(t)$. Raising $1-F(t)$ to the integer power $M$ and integrating from zero to one uses rational arithmetic only. At $(n,k,M)=(99,90,2)$ this yields~\eqref{eq:exactnumber}. The code handles $k=n+1$ as coverage one before attempting a Beta distribution. Numeric quadrature is a separate check and is never described as exact arithmetic.

\subsection{Taylor-certified orderings}
For a symmetric-chain family of $q$ permutations on $m$ points, choose distinct rational $1<t_r<2$, $\beta_r=\log t_r$, and baseline competitor gaps $-1,\ldots,-q$. Let $R_{ri}\in\{0,\ldots,m-1\}$ be point $i$'s rank in permutation $r$. The rational matrix $E_{rj}=t_r^{-j}$, $1\le j\le q$, is invertible. Solve $EV=R$ exactly, set $H=\max(1,\max_{ji}|V_{ji}|)$, $W=V/H$, and $\varepsilon=1/(12qH)$. Point $i$'s competitor logits are $-j+\varepsilon W_{ji}$, with target logit zero. They are negative. At $\beta_r$, its reciprocal probability satisfies
\[
 A_i(\beta_r)=A_0(\beta_r)+\frac{\beta_r\varepsilon}{H}R_{ri}+R_i.
\]
For $|x|\le1$, $|e^x-1-x|\le(3/2)x^2$. Since $\beta_r<1$, $|W_{ji}|\le1$, and baseline exponential factors are at most one,
\[
 \frac{|R_i|}{\beta_r\varepsilon/H}
 \le\tfrac32q\beta_r\varepsilon H\le\tfrac18.
\]
Adjacent target ranks differ by one, so their actual denominators remain separated by at least $3/4$ of the first-order step. These rational linear identities and inequalities certify the requested orderings. Evaluating exponentials to 160 decimal digits is secondary corroboration. The new files certify $(m,K)=(3,4)$ and $(4,7)$; the original failed four-class witness has different logits and remains false.

\subsection{Verification records}
The supplementary material records exact and numerical verification separately. It includes the rejected floating-point witness, the rational crossing check for the three-class boundary, both Taylor certificates, the full-set quantile boundary, all simulation outputs, and the deterministic Beta-law grid. Runtime and protocol details are kept there rather than used as evidence for a universal theorem.
\end{document}
